\chapter{Implementation}
\label{ch:impl}

This chapter presents selected implementation details that support
the architectural claims in Chapters~\ref{ch:arch}--\ref{ch:pipeline}.
Code excerpts are short and deliberately illustrative rather than
exhaustive; the full source is available in the accompanying
repository.

\section{The Event Schema}
\label{sec:event-schema}

The single most important design choice in the orchestrator is that
\code{stream\_chat} yields plain event dictionaries, not SSE strings.
This decouples the orchestrator from any particular front-end.

\begin{lstlisting}[caption={Event dictionary shapes emitted by
\code{orchestrator/loop.py::stream\_chat}.},
label={lst:events}]
{"type": "token",        "content": str}
{"type": "tool_start",   "name": str, "label": str,
                          "service": str, "arguments": dict}
{"type": "tool_progress","message": str, "percentage": float}
{"type": "tool_end",     "name": str, "label": str, "service": str,
                          "result": Any, "error": str | None,
                          "duration_ms": int}
{"type": "done",         "answer": str, "tool_calls": list[dict]}
{"type": "error",        "message": str}
\end{lstlisting}

A Streamlit consumer, a CLI, or a SSE-over-HTTP wrapper can all
consume this stream identically.

\section{Streaming Tool-Call Accumulation}
\label{sec:tool-accum}

OpenAI-compatible streaming APIs deliver tool call arguments in
partial chunks; the orchestrator accumulates them by index:

\begin{lstlisting}[caption={Argument accumulation loop (extract).},
label={lst:accum}]
if getattr(delta, "tool_calls", None):
    for tc_delta in delta.tool_calls:
        idx = tc_delta.index
        slot = pending_tool_calls.setdefault(
            idx, {"id": None, "name": "", "arguments": ""})
        if tc_delta.id:
            slot["id"] = tc_delta.id
        fn = getattr(tc_delta, "function", None)
        if fn:
            if getattr(fn, "name", None):     slot["name"]      = fn.name
            if getattr(fn, "arguments", None): slot["arguments"] += fn.arguments
        if idx not in started_tool_indices and slot["name"]:
            started_tool_indices.add(idx)
            yield {"type": "tool_start", "name": slot["name"], ...}
\end{lstlisting}

A \code{tool\_start} event is emitted as soon as the tool name is
first observed, giving the UI a chance to render an ``in-progress''
placeholder even before the argument JSON has finished streaming.

\section{Progress Bridging via \code{asyncio.Queue}}
\label{sec:queue-bridge}

Tools emit progress via \code{ctx.report\_progress(current, total,
message)}. The orchestrator installs a handler that pushes
\code{tool\_progress} events into an asyncio queue, then interleaves
the queue with the tool's completion future:

\begin{lstlisting}[caption={Interleaving tool execution and progress
events.}, label={lst:queue}]
call_task = asyncio.create_task(
    pool.call_tool(name, args, progress_handler=progress_handler))

while True:
    try:
        evt = await asyncio.wait_for(progress_queue.get(), timeout=0.05)
        yield evt
    except asyncio.TimeoutError:
        if call_task.done():
            break
while not progress_queue.empty():
    yield progress_queue.get_nowait()
\end{lstlisting}

The tight 50~ms poll interval feels responsive to the user without
starving the event loop.

\section{The MCPClientPool}
\label{sec:pool-impl}

The pool exposes a small surface but does non-trivial routing work.

\begin{lstlisting}[caption={\code{MCPClientPool.call\_tool} (extract).},
label={lst:call-tool}]
async def call_tool(self, tool_name, arguments,
                    progress_handler=None, log_handler=None):
    service = self.tool_index.get(tool_name)
    if service is None:
        raise KeyError(f"Tool not in unified registry: {tool_name}")
    actual_name = tool_name
    if "__" in tool_name and \
       tool_name.split("__", 1)[0].replace("_", "-") == service:
        actual_name = tool_name.split("__", 1)[1]
    client = self.clients.get(service)
    state = self.states[service]
    if client is None:
        state.status = "offline"
        raise RuntimeError(f"Service `{service}` is not connected")
    state.status = "processing"
    await self._notify_observers(state)
    try:
        kwargs = {}
        if progress_handler is not None:
            kwargs["progress_handler"] = progress_handler
        result = await client.call_tool(actual_name, arguments, **kwargs)
        state.status, state.last_error = "online", None
        return result
    except Exception as exc:
        state.status = "error"
        state.last_error = f"{type(exc).__name__}: {exc}"
        raise
\end{lstlisting}

The name collision handling around \code{"\_\_"} allows two servers
to expose tools with identical names without the pool rejecting the
registration.

\section{The Workspace Registry}
\label{sec:workspace-impl}

The full artefact save function is short enough to reproduce in
full:

\begin{lstlisting}[caption={\code{core/workspace.py::save\_artifact}
(complete function).}, label={lst:save-art}]
def save_artifact(*, local_path: Path, kind: str,
                  metadata: dict[str, Any] | None = None,
                  move: bool = True) -> ArtifactRecord:
    if not local_path.exists():
        raise FileNotFoundError(f"Cannot register non-existent file: {local_path}")
    category = _KIND_TO_CATEGORY.get(kind, "misc")
    dest_dir = _artifacts() / category
    dest_dir.mkdir(parents=True, exist_ok=True)
    artifact_id = str(uuid.uuid4())
    safe = _safe_stem(local_path.name)
    dest_path = dest_dir / f"{artifact_id[:8]}__{safe}"
    if move:
        shutil.move(str(local_path), dest_path)
    else:
        shutil.copy2(local_path, dest_path)
    record = ArtifactRecord(
        id=artifact_id, kind=kind, category=category,
        filename=local_path.name, path=str(dest_path.resolve()),
        size_bytes=dest_path.stat().st_size,
        mime_type=_mime_for(dest_path), metadata=metadata or {},
    )
    idx = _load_artifact_index()
    idx.append(record.to_dict())
    _write_artifact_index(idx)
    return record
\end{lstlisting}

\section{The Bake-off Worker}
\label{sec:worker-impl}

Each candidate model is trained in a worker function that captures
its own errors, so a single failing model does not abort the whole
race:

\begin{lstlisting}[caption={\code{\_train\_one} worker (extract).},
label={lst:worker}]
def _train_one(name, estimator, pre, Xtr, ytr, Xte, yte, problem, cv):
    pipe = Pipeline([("pre", pre), ("model", estimator)])
    t0 = time.time()
    scoring = "f1_weighted" if problem == "classification" else "r2"
    try:
        cv_scores = cross_val_score(pipe, Xtr, ytr, cv=cv,
                                    scoring=scoring, n_jobs=-1)
        pipe.fit(Xtr, ytr)
        preds = pipe.predict(Xte)
        metrics = _metrics(problem, ytr, yte, preds)
        return {"model": name, "cv_mean": float(cv_scores.mean()),
                "cv_std": float(cv_scores.std()), "metrics": metrics,
                "train_seconds": round(time.time()-t0, 3),
                "pipeline": pipe, "error": None}
    except Exception as exc:
        return {"model": name, "cv_mean": None, "cv_std": None,
                "metrics": {}, "train_seconds": round(time.time()-t0, 3),
                "pipeline": None,
                "error": f"{type(exc).__name__}: {exc}"}
\end{lstlisting}

\section{Streamlit \code{PoolRunner}}
\label{sec:st-runner}

Streamlit's model re-runs the script top-to-bottom on every widget
event. To keep the MCP pool alive across those re-runs, it is cached
via \code{st.cache\_resource} inside a \code{PoolRunner} that owns a
dedicated event loop:

\begin{lstlisting}[caption={\code{PoolRunner} (extract).},
label={lst:runner}]
class PoolRunner:
    def __init__(self):
        self.loop = asyncio.new_event_loop()
        self.thread = threading.Thread(target=self._run, daemon=True)
        self.thread.start()
        self.pool = MCPClientPool()
        register_default_servers(self.pool)
        asyncio.run_coroutine_threadsafe(self.pool.start(), self.loop).result()

    def _run(self):
        asyncio.set_event_loop(self.loop)
        self.loop.run_forever()

    def stream_events(self, query, active_file, history):
        q: queue.Queue = queue.Queue()
        async def _drive():
            try:
                async for evt in stream_chat(
                    query=query, active_file=active_file,
                    history=history, pool=self.pool):
                    q.put(evt)
            finally:
                q.put(None)
        asyncio.run_coroutine_threadsafe(_drive(), self.loop)
        return q
\end{lstlisting}

The thread-safe queue is the bridge: the async producer runs in the
background loop, while the Streamlit script consumes events
synchronously and updates placeholders in place.

\section{Error Handling Discipline}
\label{sec:errors-impl}

Errors are surfaced consistently through the event pipeline. Tools
that hit a domain error (e.g.\ ``target column not found'') return a
result dictionary of the form \code{\{"error": "\dots"\}}; tools
that raise a Python exception have the exception caught by the
orchestrator, converted into a synthetic result dictionary with
\code{\{"error": "TypeName: message"\}}, and surfaced as a
\code{tool\_end} event with a non-null \code{error} field. This
makes the LLM's next reasoning step deterministic: it sees the
error string in the tool result and can decide whether to retry, ask
the user for clarification, or continue with a different plan.

\section{Test and Smoke Verification}
\label{sec:tests-impl}

The repository currently contains no automated test suite;
verification is done through a smoke script that:

\begin{enumerate}[leftmargin=1.4em]
  \item Imports and instantiates the pool.
  \item Calls \tool{list\_uploaded\_files} and confirms an empty
        registry starts empty.
  \item Uploads a small synthetic CSV via
        \code{save\_uploaded\_dataset} and calls
        \tool{ingest\_dataset} to confirm the round trip.
\end{enumerate}

The absence of automated tests is a known limitation and is discussed
in Chapter~\ref{ch:limitations}.
